\documentclass[../../../../main.tex]{subfiles}
\begin{document}

\begin{flushright}
\footnotesize\textit{【自動文字起こし・要確認】}
\end{flushright}

[解] $A, B, C$の勝つ確率は右図のようである。\\
\begin{minipage}{0.4\textwidth}
\begin{tabular}{|c|c|c|c|}
\hline
相手 & A & B & C \\
\hline
A & \diagbox{}{} & $p$ & $q$ \\
\hline
B & $1-p$ & \diagbox{}{} & $\frac{1}{2}$ \\
\hline
C & $1-q$ & $\frac{1}{2}$ & \diagbox{}{} \\
\hline
\end{tabular}
\end{minipage}

(1) もとめる確率を$P_a$とする。右図から、
\[ P_a = pq + p(1-q) \cdot \frac{1}{2} \cdot P_a \]
\[ \therefore P_a = \frac{2pq}{2-p(1-q)} \]

\begin{tikzpicture}[x=0.5cm, y=-0.5cm, baseline=(current bounding box.north)]
\node at (0,0) {A};
\node at (2,0) {B};
\node at (4,0) {C};
\node at (6,0) {B};
\node at (8,0) {A};
\draw (0,0.5) -- (0,1) -- (2,1) -- (2,0.5);
\draw (1,1) -- (1,1.5) -- (4,1.5) -- (4,0.5);
\draw (2.5,1.5) -- (2.5,2) -- (6,2) -- (6,0.5);
\draw (4.25,2) -- (4.25,2.5) -- (8,2.5) -- (8,0.5);
\draw[dashed] (6.125,2.5) -- (6.125,3.5);
\end{tikzpicture}

(2) もとめる確率を$P_b$とする。右図から、
\[ P_b = (1-p) \frac{1}{2} \cdot P_c \]
$P_c$は(1)で$p$と$q$を入れかえたものだから、
\[ P_b = \frac{pq(1-p)}{2-q(1-p)} \]

\begin{tikzpicture}[x=0.5cm, y=-0.5cm, baseline=(current bounding box.north)]
\node at (0,0) {A};
\node at (2,0) {B};
\node at (4,0) {C};
\node at (6,0) {A};
\draw (0,0.5) -- (0,1) -- (2,1) -- (2,0.5);
\draw (1,1) -- (1,1.5) -- (4,1.5) -- (4,0.5);
\draw (2.5,1.5) -- (2.5,2) -- (6,2) -- (6,0.5);
\draw[dashed] (4.25,2) -- (4.25,3);
\node[right] at (4.5, 3) {これ以降Aが優勝する確率が $P_c$ となる。};
\end{tikzpicture}

(3) もとめる確率を$R$とする。1回目でどちらが勝つかで場合分けして、
\[ R = \frac{1}{2} (P_a + P_c) = pq \left\{ \frac{1}{2-p(1-q)} + \frac{1}{2-q(1-p)} \right\} \]

\end{document}
